\subsection{Wissenschaftlicher Ansatz}
\label{subsec:01-05-01_sci-approach}

Die Arbeit folgt einem überwiegend qualitativen, fallstudienbasierten Forschungsansatz und untersucht den Wiki-Auftritt einer organisationalen Transformationseinheit als \enquote{Single Case} im Sinne der Fallstudiensystematik nach Yin \cite{404:Yin2018}. Ziel ist es, die Wechselwirkungen zwischen \Gls{ia}, \Gls{ux}-Strategien und der Akzeptanz interner Wissensplattformen im realen Nutzungskontext zu verstehen, anstatt generalisierbare Ursache-Wirkungs-Zusammenhänge experimentell zu testen. Der zugrunde liegende Erkenntnisprozess orientiert sich an einem deduktiv-induktiven Vorgehen, das theoretische Konzepte mit empirisch gewonnenen Einsichten verbindet. \autoref{fig:g-05_methodology-of-inductive-approach} veranschaulicht die methodische Logik dieses Vorgehens.

\image{H}{0.65}{g-05_methodology-of-inductive-approach}{Methodisches Vorgehen eines deduktiv-induktiven Forschungsansatzes: Verbindung theoretischer Konzepte mit empirischer Datenerhebung und -analyse (\gls{iAa} \cite{200:ScribbrInductiveDeductive2023})}

Wie die Abbildung zeigt, werden bestehende theoretische Annahmen aus der \Gls{ux}- und \Gls{ia}-Forschung als analytischer Rahmen genutzt, zugleich aber durch qualitative Datenerhebung und induktive Musterbildung weiterentwickelt. Methodisch kombiniert die Arbeit eine systematische Literaturrecherche, ein leitfadengestütztes Experteninterview, die konzeptionelle Weiterentwicklung eines prototypischen Wiki-Auftritts sowie darauf aufbauende Usability-Tests und Feedbackinterviews.

Die Literaturrecherche bildet die Grundlage der theoretischen Fundierung in \hyperref[ch:02_theoretical-foundations]{Kapitel~2} sowie der Einordnung der Ergebnisse in \hyperref[ch:06_discussion]{Kapitel~6}. Zentrale Quellen sind peer-reviewte Publikationen aus \emph{IEEE Xplore} zu \Gls{ux}, Human-Centred Design und Technologieakzeptanz \cite{902:IEEEXplore}. Ergänzend wird \emph{Google Scholar} zur Erschließung von Zitationsnetzwerken und aktuellen Veröffentlichungen genutzt \cite{903:GoogleScholar}. Für praxisnahe Themen, etwa Intranet-Informationsarchitekturen oder Usability-Best-Practices, erfolgt zusätzlich eine gezielte Webrecherche unter Berücksichtigung renommierter Institutionen und Fachorganisationen.

Die Suchstrategie umfasst drei Themenfelder: (A) \Gls{ux}, \Gls{ia} und Human-Centred Design, (B) Technologieakzeptanz sowie (C) interne Wissensplattformen und Intranets. Die identifizierten Quellen werden systematisch erfasst und anhand eines sechsstufigen Relevanzschemas bewertet, das u.a. Peer-Review-Status, Aktualität, methodische Transparenz und Zitierhäufigkeit berücksichtigt \cite{200:ScribbrInductiveDeductive2023}. Hoch relevante Quellen bilden den Kern der theoretischen Argumentation.

Zur Erhebung kontextbezogenen Expertenwissens wird ein leitfadengestütztes, halbstrukturiertes Interview nach Kvale und Brinkmann durchgeführt \cite{405:KvaleBrinkmann2015}. Der Interviewleitfaden fokussiert Nutzungsbarrieren, Gestaltungsprinzipien und didaktische Aspekte der Wissensaufbereitung in Wikis. Die Ergebnisse fließen in die Voranalyse, die theoretische Rahmung sowie die Diskussion ein und unterstützen insbesondere die Beantwortung von RQ-1 und RQ-3.

Die Weiterentwicklung des Wiki-Auftritts basiert auf einem bestehenden Proof-of-Concept. Aufbauend auf theoretischen Gestaltungsprinzipien und den Interviewergebnissen werden zentrale \Gls{ia}-Entscheidungen (z.B. Navigationsstruktur, Benennungskonventionen) sowie \Gls{ux}-Muster (z.B. visuelle Orientierungselemente) prototypisch umgesetzt. Dieses Artefakt dient als Untersuchungsgegenstand der Evaluation und verknüpft den gestaltungsorientierten Ansatz mit empirischer Analyse im Sinne des Design-Science-Research-Paradigmas \cite{304:FraunhoferDSR2015}.

Zur Evaluation werden formative Usability-Tests mit zehn Personen durchgeführt, die unterschiedliche Rollen und Zugehörigkeiten abdecken. Die Auswahl erfolgt auf Basis zuvor entwickelter Personas. In zwei Testdurchläufen bearbeiten die Teilnehmenden definierte Aufgaben im bestehenden und im neuen Wiki; erfasst werden Erfolgsraten, Bearbeitungszeiten, Navigationsstrategien und Fehler \cite{104:Nielsen2024FiveUsers}. Ergänzend werden kurze Feedbackinterviews zur subjektiven Wahrnehmung von Verständlichkeit, Orientierung und Akzeptanz geführt \cite{405:KvaleBrinkmann2015}. Die qualitativen Befunde bilden die Hauptgrundlage für RQ-2 und RQ-3, während die quantitativen Metriken explorativ interpretiert werden \cite{404:Yin2018,601:Schadewitz2007}.

Die Integration der Daten erfolgt in einem triangulativen, deduktiv-induktiven Analyseverfahren \cite{404:Yin2018}. Theoretisch abgeleitete Kategorien dienen als Analysegerüst und werden im Auswertungsprozess induktiv ergänzt \cite{601:Schadewitz2007,200:ScribbrInductiveDeductive2023}. Die Ergebnisse werden in \hyperref[ch:06_discussion]{Kapitel~6} mit den Forschungsfragen verknüpft und bilden die Grundlage für die Handlungsempfehlungen in \hyperref[ch:07_recommendations]{Kapitel~7}.